\section{Scale gluing at the critical distribution exponent}
\label{sec:v91-gluing}
The critical surface estimate has more information than a bound on
one fixed window. Its constant is linear in the length of a window
matched to the exit width. A finite partition retains this gain.

\begin{lemma}[Gluing with the local scale retained]
\label{lem:v91-gluing}
Let $q>1$, $\beta=1-1/q$, and partition an interval $J$ into $N$
measurable pieces, each contained in an interval on which
$\mathcal N_q(f;\cdot)\le A$. Then
\begin{equation}\label{eq:v91-gluing}
                 \mathcal N_q(f;J)\le N^{1/q}A.
\end{equation}
Consequently, for source families $f_{m,r}$, if
\[
 \limsup_m\sup_{R,n,k}\sup_{|I|=r}
                 \mathcal N_q(f_{m,r};I)\le C r^\zeta,
\]
then, for every fixed $L,r>0$,
\begin{equation}\label{eq:v91-local-global}
 \limsup_m\sup_{R,n,k}\sup_{|J|=L}
 \mathcal N_q(f_{m,r};J)
       \le C(L+r)^{1/q}r^{\zeta-1/q}.
\end{equation}
The interval partition is chosen before the collision limit.
\end{lemma}
\begin{proof}
For $E\subset J$ put $v_i=|E\cap J_i|$. At each finite count,
\[
 \int_E|f|\le A\sum_{i=1}^N v_i^\beta
             \le A N^{1-\beta}|E|^\beta.
\]
This is H\"older's inequality on a finite counting measure. Divide
by $|E|^\beta$ and take the measurable-set supremum. For the second
assertion partition into $N=\lceil L/r\rceil$ pieces of length at
most $r$ and enlarge the last piece for its upper estimate. Now
$N\le(L+r)/r$. The first inequality is valid before any parameter
supremum, so it also holds for moving interval locations. For fixed
$L,r$, $N$ is independent of $m$. No infinite sum of collision
limsups occurs.
\end{proof}

\begin{theorem}[Complete critical surface decay]
\label{thm:v91-surface}
For every fixed $L>0$ and sufficiently small fixed $r>0$,
\begin{equation}\label{eq:v91-surface}
 \limsup_m\sup_{R,n,k}\sup_{|J|=L}
 \mathcal N_Q(m^2\overline{\mathcal F}_r;J)
                  \le C(L+r)^{31/32}r^{1/32}.
\end{equation}
Every positive residence length, every incoming endpoint type and
all one-sided incident physical components are included.
\end{theorem}
\begin{proof}
The retained critical surface inequality, evaluated with $|I|=r$,
gives \eqref{eq:v91-matched-input}. Apply
Lemma~\ref{lem:v91-gluing} with $q=Q$, $\zeta=1$ and
$f_{m,r}=m^2\overline{\mathcal F}_r$. The local upper bound already
sums all residence scales at finite collision count. The new
partition concerns roof intervals only; it neither partitions the
physical source into a selected word nor introduces an additional
residence cutoff.
\end{proof}

\begin{lemma}[The fractional potential at the same scale]
\label{lem:v91-potential}
For $v=v^{1/4}$ and $\mathsf P_r v(t)=\int_0^r s^{-3/8}v(t+s)\,ds$,
\begin{equation}\label{eq:v91-potential}
 \limsup_m\sup_{R,n,k}\sup_{|J|=L}
 \|m^2\mathsf P_r v\|_{L^Q(J)}
                   \le C(L+r)^{31/32}r^{5/8}.
\end{equation}
\end{lemma}
\begin{proof}
On an interval $I$ of length $r$, Minkowski's inequality gives
\[
 \|m^2\mathsf P_r v\|_{L^Q(I)}
 \le (1-3Q/8)^{-1/Q}r^{1/Q-3/8}
                          m^2\int_{I+(0,r)}v.
\]
Here $3Q/8=12/31<1$. The fixed-width source estimate bounds the
ordered mass on the right by $Cr$, uniformly in its location.
Thus the local $L^Q$ norm is at most $Cr^{1/Q+5/8}$ in the
collision limsup. Partition $J$ into at most $(L+r)/r$ such
intervals, sum the $Q$th powers at finite count and take the limit.
This proves \eqref{eq:v91-potential}. The width $r$ stays fixed
through the collision and auxiliary-band limits.
\end{proof}

\begin{theorem}[The critical raw limit without logarithmic loss]
\label{thm:v91-raw}
Under the retained inputs, for every fixed $M,L>0$,
\begin{equation}\label{eq:v91-raw}
 \sup_{R,n,k}\sup_{|J|=L}
 \mathcal N_Q\bigl((m^2p-\mathcal L)\one_{\mathcal C(M)};J\bigr)
                                      \longrightarrow0.
\end{equation}
More precisely, for every fixed $\varepsilon>0$,
\begin{equation}\label{eq:v91-bad-source}
 \limsup_m\sup_{R,n,k}\sup_{|J|=L}
       \mathcal N_Q(m^2b^\varepsilon;J)
                         \le C L^{31/32}\varepsilon^{1/16}.
\end{equation}
Neither conclusion asserts strong $L^Q$ or essential-supremum
convergence.
\end{theorem}
\begin{proof}
Apply $\mathcal N_Q$ to the positive source comparison
\eqref{eq:v91-source-comparison}. Its triangle inequality follows
directly from its definition. For an essentially bounded function,
$\mathcal N_Q(f;J)\le L^{1/Q}\|f\|_\infty$; also
$\mathcal N_Q(f;J)\le\|f\|_{L^Q(J)}$. The fixed-width estimates,
Theorem~\ref{thm:v91-surface} and Lemma~\ref{lem:v91-potential}
therefore give
\begin{align}
 \limsup_m\sup_{R,n,k}\sup_{|J|=L}
       \mathcal N_Q(m^2b^\varepsilon;J)
 \le C_\Theta(L+r)^{31/32}
  \{\varepsilon^{1/16}+r^{1/32}+r^{5/8}
                       +(1+\varepsilon^{-1})\sqrt r\}.
 \label{eq:v91-source-budget}
\end{align}
All norms are applied to the complete source before limits.
Let $r\downarrow0$ only after the collision limit and the permitted
auxiliary-band enlargement, with $\varepsilon$ fixed. This proves
\eqref{eq:v91-bad-source}.

Use the original positive scalar remainder with
$\varepsilon(B)=A_0B^{-1/12}$. Restricting to $\mathcal C(M)$
can only decrease a positive norm. The remainder has critical norm
at most $C_M L^{31/32}B^{-1/192}$ in the ordered limit, as does
the bad-source term. Send $B\to\infty$ last to prove
\eqref{eq:v91-raw}. Equivalently choose $r(B)=c\varepsilon(B)^4$
in \eqref{eq:v91-source-budget}, fixed before $m\to\infty$;
all its explicit errors are $O(\varepsilon(B)^{1/16})$.
No frequency band depends on the collision count.
\end{proof}

\begin{corollary}[The full central window]
\label{cor:v91-diffusive}
For $H=[-M,M]$ put
\[
 \widetilde f_{m,R,n,k}(u)=
 (m^2p-\mathcal L)(m\bar\tau_R+\sqrt m\,u)
 \one_{\mathcal C(M)}(m\bar\tau_R+\sqrt m\,u).
\]
Then $\sup_{R,n,k}\mathcal N_Q(\widetilde f_{m,R,n,k};H)\to0$.
In roof coordinates the measure is $dt/\sqrt m$, not $dt$.
\end{corollary}
\begin{proof}
Partition the central roof window into at most $2M\sqrt m+2$
intervals of length at most one. If $d_m\to0$ is the uniform
local norm from \eqref{eq:v91-raw}, the glued norm is at most
$d_m(2M\sqrt m+2)^{1/Q}$. Changing to $u$ multiplies the norm
by $m^{-1/(2Q)}=m^{-31/64}$. The product tends to zero. The
growing number of intervals is paid explicitly, not treated as fixed.
\end{proof}
